\section*{Chapitre \ref{ch:dv:implied-volatility-and-its-surface} --- Volatilité implicite et sa surface}

\begin{solution}{exo:dv:implied-volatility-and-its-surface:1}
$k=\ln(90/100.37)=-0.1091$~; $w=0.2047^2\times91/365=0.0104$.
\end{solution}

\begin{solution}{exo:dv:implied-volatility-and-its-surface:2}
$20.47-14.19=6.27$ points de volatilité. À un an, les mêmes strikes sont plus proches de la monnaie en
écarts-types ($\ln(K/F)/\sqrt T$ est divisé par deux), et le skew par unité de $k$ est lui-même
plus faible, décroissant à peu près comme $T^{-1/2}$~: le skew 90--110 à un an vaut 3.3 points.
\end{solution}

\begin{solution}{exo:dv:implied-volatility-and-its-surface:3}
$w$ vaut $0.2^2\times0.25=0.0100$ à trois mois et $0.13^2\times0.5=0.0085$ à six mois~: la
variance totale baisse, c'est un \omterm{def:dv:implied-volatility-and-its-surface:static}{arbitrage calendaire}. Vendre le straddle à trois mois (ou le call) et acheter
celui à six mois, à même moneyness forward~: l'option de plus longue échéance vaut au moins
celle de courte échéance à la première échéance.
\end{solution}

\begin{solution}{exo:dv:implied-volatility-and-its-surface:4}
Sans le skew~: $P\Phi(d_2)=0.4980$. La correction vaut $-\mathcal V\,\partial\sigma/\partial K
=19.77\times0.0031=0.0604$ (par unité de volatilité)~: 0.5584 avec le skew.
\end{solution}

\begin{solution}{exo:dv:implied-volatility-and-its-surface:5}
\omterm{def:dv:implied-volatility-and-its-surface:sticky}{Sticky strike}~: le delta de Black--Scholes, 0.533. \omterm{def:dv:implied-volatility-and-its-surface:sticky}{Sticky delta}~: $0.533+19.77\times0.00306=0.593$.
Six actions pour 100 calls.
\end{solution}

\begin{solution}{exo:dv:implied-volatility-and-its-surface:6}
La distribution risque-neutre d'un indice ajoute la tendance de ses composants à chuter ensemble~: dans un krach
les corrélations montent, si bien que la queue gauche de l'indice est plus épaisse que celle de chaque
composant rapporté à sa propre volatilité. Les distributions risque-neutres individuelles sont bien moins asymétriques que
celle de l'indice~; l'écart correspond au skew de corrélation du chapitre 17.
\end{solution}

\begin{solution}{exo:dv:implied-volatility-and-its-surface:7}
Six violations d'\omterm{def:dv:implied-volatility-and-its-surface:static}{arbitrage butterfly}, la pire étant $g=-2.06$ pour $k=-0.1$. Une erreur à la baisse de 6 points sur
la cotation à six mois à la monnaie laisse $w$ croissante~; 6.5 points (le premier pas
d'un demi-point qui le fait) crée une violation d'\omterm{def:dv:implied-volatility-and-its-surface:static}{arbitrage calendaire}~; 7 points en créent une, 8 points en créent trois.
\end{solution}

\begin{solution}{exo:dv:implied-volatility-and-its-surface:8}
La volatilité implicite à un strike est un prix, et non la volatilité attendue si l'indice atteint ce
strike. Le put 80 à 27 résume en un seul nombre la probabilité et l'ampleur d'une baisse, y compris
une prime couvrant l'assurance et les sauts, que nulle volatilité de diffusion à 80 ne décrit. Ce que le
marché anticiperait sachant une baisse relève de la volatilité locale du chapitre 9 ou de la
dynamique du modèle, et non de la volatilité implicite à ce strike.
\end{solution}

\begin{solution}{pb:dv:implied-volatility-and-its-surface:1}
\textbf{1.} Forward 100.37~; facteur d'actualisation 0.99255.
\textbf{2.} 16.36~\% et 26.78~\%.
\textbf{3.} $\ln(80/100.37)=-0.2269$.
\textbf{4.} 0.219.
\textbf{5.} 0.006.
\textbf{6.} 0.31~\%.
\textbf{7.} 1.89~\%.
\textbf{8.} La probabilité vaut l'opposé de la dérivée du put par rapport au strike, divisée par le
facteur d'actualisation, et la volatilité du put varie avec le strike~: le skew ajoute $\mathcal V\,
\partial\sigma/\partial K$, terme que ignore $\Phi(-d_2)$ à volatilité fixe.
\textbf{9.} C'est la valeur à terme d'un put digital de strike 80, et la limite d'un put spread
$(P(80+h)-P(80-h))/(2h)$ divisée par le facteur d'actualisation.
\textbf{10.} Du centre et de la queue droite vers la queue gauche (et un peu vers la région
juste au-dessus du forward, où le pic monte).
\textbf{11.} Non~: c'est un prix, il s'agit de la probabilité risque-neutre, celle qui fait des prix des espérances,
et qui inclut la prime que les investisseurs paient pour se protéger de la baisse.
\textbf{12.} La prime de risque de l'assurance contre le krach, la demande des hedgers, et le coût
de portage du risque supporté par les dealers.
\textbf{13.} Le put profite de la baisse dans les deux régimes~; sous \omterm{def:dv:implied-volatility-and-its-surface:sticky}{sticky strike}, sa volatilité reste
inchangée à 26.8, sous \omterm{def:dv:implied-volatility-and-its-surface:sticky}{sticky delta} elle baisse, puisque le put est désormais plus proche de la monnaie
(23.8 au nouveau moneyness)~: le \omterm{def:dv:implied-volatility-and-its-surface:sticky}{sticky strike} donne le gain le plus important.
\textbf{14.} Le krach a montré qu'une baisse d'un cinquième en deux jours était possible~; les acheteurs de
protection et les vendeurs ayant perdu ont réévalué la queue gauche, et continuent de le faire.
\textbf{15.} Une baisse est soudaine~: un saut ou une poussée de volatilité pèse davantage dans la variance
à courte échéance qu'à longue échéance, si bien que les skews de courte échéance sont plus pentus (chapitres 12--13).
\textbf{16.} Que le marché facture environ six fois la probabilité à volatilité plate, et que
le prix soit de 0.22 pour 100 de notionnel par trimestre, soit environ 0.9~\% par an si on le reconduit.
\textbf{17.} Un put spread (achat 80, vente 70) ou un collar~: moins cher, mais la protection s'arrête sous
70, ou l'on abandonne la hausse.
\textbf{18.} La surface avec skew, celle où la protection peut réellement s'acheter~; la
surface plate sous-estime la queue d'un facteur six.
\textbf{19.} 0.31~\% à plat, 1.89~\% avec skew~: six fois plus.
\textbf{20.} Le prix que le marché attribue à la queue gauche de la distribution, au-delà de ce qu'une seule
volatilité peut décrire.
\end{solution}

\begin{solution}{iq:dv:implied-volatility-and-its-surface:1}
Les investisseurs achètent des puts sur indice comme assurance, les krachs vont à la baisse et sont
soudains, et la volatilité monte quand les cours baissent~; avant 1987, le marché évaluait des queues
presque lognormales, et le krach a montré qu'une baisse d'un cinquième en un ou deux jours était possible. Les vendeurs de puts le facturent depuis.
\iqlookfor{demande d'assurance, sauts, effet de levier et rupture historique.}
\end{solution}

\begin{solution}{iq:dv:implied-volatility-and-its-surface:2}
\omterm{def:dv:implied-volatility-and-its-surface:density}{Formule de Breeden--Litzenberger}~: $q(K)=\partial^2C/\partial K^2/P(0,T)$. En pratique, ajuster un
smile de volatilité sans arbitrage (les cotations brutes sont trop bruitées pour être dérivées deux fois), puis dériver
les prix des calls lissés~; vérifier que la densité s'intègre à un et reproduit le forward.
\iqlookfor{la formule et la nécessité de lisser d'abord.}
\end{solution}

\begin{solution}{iq:dv:implied-volatility-and-its-surface:3}
\omterm{def:dv:implied-volatility-and-its-surface:sticky}{Sticky strike}~: chaque strike garde sa volatilité, si bien que la nouvelle volatilité à la monnaie se lit
plus haut sur le skew (18.1 contre 16.4 dans le chapitre) et le delta est le delta de Black--Scholes.
\omterm{def:dv:implied-volatility-and-its-surface:sticky}{Sticky delta}~: le smile se déplace avec le forward, la volatilité à la monnaie est inchangée, et
le delta gagne le \omterm{def:dv:greeks-and-the-hedging-pnl:greeks}{vega} multiplié par la pente, terme d'autant plus fort que le skew est négatif.
\iqlookfor{le changement de niveau et le changement de couverture.}
\end{solution}

\begin{solution}{iq:dv:implied-volatility-and-its-surface:4}
Les calls sont décroissants en strike et ont une pente supérieure à $-P(0,T)$~; convexes en strike (butterflies
non négatifs), soit de façon équivalente $g(k)\ge0$ en variance totale~; la variance totale est non décroissante
en échéance à moneyness forward fixe. Vérifier chaque tranche sur une grille fine de $k$ et chaque paire de
tranches consécutives~; avec cours acheteur et cours vendeur, vérifier aux prix exécutables.
\iqlookfor{les trois conditions et leur forme en variance totale.}
\end{solution}

\begin{solution}{iq:dv:implied-volatility-and-its-surface:5}
Elle utilise une queue lognormale, en ignorant le skew. La probabilité est la dérivée en strike du
put, qui inclut le terme de skew~; sur une surface d'indice typique, l'erreur à 20~\% en dessous est
d'un facteur cinq ou plus (six dans le chapitre).
\iqlookfor{le skew dans la dérivée, et l'ampleur de l'erreur.}
\end{solution}

\begin{solution}{iq:dv:implied-volatility-and-its-surface:6}
Travailler en variance totale et en \omterm{def:dv:implied-volatility-and-its-surface:surface}{log-moneyness}~; à échéance fixée, utiliser une paramétrisation lisse ou une
spline en $k$ qui respecte $g\ge0$ (SVI, chapitre 8), avec précaution aux extrémités~; entre les échéances
interpoler $w$ linéairement dans le temps à $k$ fixé, ce qui préserve la monotonie calendaire~; traiter
les dividendes et les événements dans la coordonnée de temps~; vérifier le résultat.
\iqlookfor{variance totale, interpolation monotone dans le temps, et un contrôle.}
\end{solution}
